\ \ \\[2cm]
A crescente integração de animais de estimação como membros da família, somada à rotina atribulada dos tutores, impulsiona a busca por soluções tecnológicas que garantam o bem-estar animal. Este trabalho apresenta a concepção, validação e evolução de um sistema de baixo custo baseado em Internet das Coisas (IoT) para o monitoramento e automação de ambientes para animais de estimação. O protótipo, validado em simulação de alta fidelidade (TCC 1), utiliza um microcontrolador ESP32 para coletar dados de múltiplos sensores e controlar atuadores.

Como evolução (TCC 2), e em resposta a uma análise de requisitos de Engenharia de Software, o trabalho implementa duas extensões avançadas. Primeiramente, um módulo de saúde preditiva, que aplica técnicas de \textit{Machine Learning} (Isolation Forest) sobre um \textit{dataset} controlado para detectar anomalias nos padrões de consumo, validando a inteligência do sistema. Em segundo lugar, a arquitetura é migrada para uma plataforma de nuvem escalável e segura (AWS IoT Core e Lambda), permitindo a integração e o controle robusto por comandos de voz através de uma \textit{Skill} da Amazon Alexa.

A orquestração da lógica, a visualização de dados em \textit{dashboard} e o envio de notificações são gerenciados pela ferramenta de programação visual Node-RED. O objetivo é fornecer uma solução integrada que assegure a alimentação, hidratação e conforto térmico do animal, permitindo o controle remoto e por voz, e fornecendo \textit{insights} preditivos sobre o bem-estar do pet.

\vspace{1cm}
Palavras-chave: Internet das Coisas; Engenharia de Software; Detecção de Anomalias; AWS; Alexa; ESP32; Node-RED; Bem-Estar Animal.